% !TEX root = ../../../main.tex
% Block card: detection.cwt_ridges -- "CWT ridges" (Part I, the blocks chapter).
\cardfield{What it does}
Finds bumps (spikes) and dips (troughs) of every duration in one pass, from seconds to hours,
and says how long each one is.

\cardfield{When to use it}
When the events you care about come in very different lengths and you do not want to choose one
window size in advance. It also copes with a slowly drifting baseline without a separate
detrending step. Put \blockref{interrogation.event_shape} after it when you need each event's
onset and recovery: that pair is the \texttt{cwt\_events} template.

\cardfield{What goes in, what comes out}
In goes a span of one channel. Out comes one stretch per event, centred on it, labelled
\texttt{spike} or \texttt{drop}, scored by how far it stands out from the noise at its own
duration. Each event also carries its own scale and an estimated duration.

\begin{plainterms}[Spikes and troughs at every scale]
Imagine a rubber stamp shaped like a Mexican hat: a hump with a shallow dip on each side.
Press it along the trace at many sizes, from small to large. A bump in the signal matches the
stamp best at the size closest to its own width; a dip matches the stamp turned upside down.
Following the best match across neighbouring sizes traces a ``ridge''; a long, strong ridge is an
event, and the size at which it is strongest tells you how long the event lasts. Doing this for
the hat both ways up and at every size is what ``spikes and troughs at every scale'' means.
\end{plainterms}

\cardfield{Settings that matter}
\emph{min\_snr} (3) first: how many times the local noise at an event's own duration it must
reach. \emph{min\_scale\_s} and \emph{max\_scale\_s} set the shortest and longest events it
looks for (an event lasts about 1.67 times its scale). \emph{polarity} chooses spikes, drops or
both in the output; both are always traced.

\cardfield{Watch out for}
Keep \emph{max\_scale\_s} below the usual gap between events: at larger sizes the hat covers
several events and can report their arrangement as one slow event. Events near the ends of the
span are reported but flagged; so are events that may be one corner of a sudden step. For a
lopsided event such as a sharkfin the centre sits towards the slow side.
